% Generated draft: requires mathematical and source review.
\documentclass[11pt]{article}
\usepackage[margin=25mm]{geometry}
\usepackage{fontspec}
\setmainfont{Latin Modern Roman}
\usepackage{amsmath,amssymb,mathtools,unicode-math}
\setmathfont{Latin Modern Math}
\usepackage{enumitem,needspace,xurl,hyperref}
\setlength{\parindent}{0pt}
\setlength{\parskip}{5pt}
\setlength{\emergencystretch}{3em}
\newcommand{\N}{\mathbb{N}}
\newcommand{\Z}{\mathbb{Z}}
\newcommand{\Q}{\mathbb{Q}}
\newcommand{\R}{\mathbb{R}}
\newcommand{\C}{\mathbb{C}}
\newcommand{\F}{\mathbb{F}}
\newcommand{\A}{\mathbb{A}}
\newcommand{\PP}{\mathbb{P}}
\newcommand{\E}{\mathbb{E}}
\newcommand{\eps}{\varepsilon}
\newcommand{\dd}{\,\mathrm{d}}
\begin{document}
\section*{UnsolvedMath ID: 701027}
\textbf{Source problem code:} catalogue entry
\subsection{Definitions and mathematical statement}
A pro-$p$ group, for a prime $p$, is a topological group isomorphic to an inverse limit of finite $p$-groups. A commutative profinite ring with identity is an inverse limit of finite associative commutative rings with identity. Equip $\operatorname{GL}_n(R)$ with its natural topology inherited from $R^{n^2}$.

Let $F$ be a free pro-$p$ group. For finite rank $d$, this means the pro-$p$ completion of the abstract free group on $d$ generators. The supplied source does not specify the rank of $F$, whether infinite rank is allowed, or whether $p$ is fixed or universally quantified. These qualifications remain unresolved; in particular, a restriction to nonabelian free pro-$p$ groups is not explicitly stated.

Under the usual interpretation of representation as a faithful continuous representation of finite dimension, the question is whether there exist an integer $n\geq 1$, an associative commutative profinite ring $R$ with identity, and a continuous injective homomorphism
\[
\rho:F\longrightarrow\operatorname{GL}_n(R).
\]
Such a map identifies $F$ topologically with a closed subgroup of $\operatorname{GL}_n(R)$. Continuity, faithfulness, and finite dimension are the interpretation used here; the source does not explicitly state these conventions.

For the accompanying formulation, call a pro-$p$ group $G$ linear if it admits such an embedding. A nontrivial pro-$p$ identity for $G$ is an element $w\neq 1$ of a free pro-$p$ group $F_m^{(p)}$ on finitely many generators $x_1,\ldots,x_m$, with $m\geq 1$, such that
\[
\varphi_{g_1,\ldots,g_m}(w)=1_G
\qquad\text{for every }(g_1,\ldots,g_m)\in G^m,
\]
where $\varphi_{g_1,\ldots,g_m}:F_m^{(p)}\to G$ is the unique continuous homomorphism sending $x_i$ to $g_i$.

The source states that a negative answer to the representation question is equivalent to every linear pro-$p$ group satisfying a nontrivial pro-$p$ identity. In this formulation, the identity may depend on the group. This equivalence is retained as a source assertion; its precise scope depends on the unresolved rank and representation conventions above.

The supplied record further reports that the identity assertion holds in dimension $2$ for $p\neq 2$, attributing this to A. N. Zubkov (1987). It also reports that a $2$-dimensional linear pro-$2$ group in characteristic $2$ satisfies a nontrivial pro-$2$ identity, attributing this to D. E.-C. Ben-Ezra and E. Zelmanov (2021). With the representation convention above, the latter characteristic condition is $\operatorname{char}(R)=2$. These are supplied literature claims, not independently verified conclusions of this draft.
\subsection{Short English statement}
The question asks whether a free pro-$p$ group has a faithful continuous representation by finite matrices over a commutative profinite ring with identity. The intended rank and prime quantifiers require clarification. The source gives an equivalent identity formulation and reports two dimension-$2$ cases of that formulation.
\subsection{Sources}
\begin{itemize}
\item Local UnsolvedMath registry, permanent ID $701027$, exact source code \texttt{catalogue entry}, title \emph{Kourovka Notebook Problem 14.42}. The supplied record labels the problem \texttt{open}; this is a registry classification, not an independently verified status assessment. Record updated: \texttt{2026-10-16T14:00:00Z}.
\item \emph{The Kourovka Notebook (Unsolved Problems in Group Theory)}, original Problem 14.42, 14th issue (1999); supplied edition reference: 21st issue edition (2026), \texttt{arXiv:1401.0300v48}. Proposed by A. N. Zubkov and V. N. Remeslennikov. The registry describes the collection as edited by E. I. Khukhro and V. D. Mazurov.
\item A. N. Zubkov, \emph{Siberian Math. J.}, volume 28, no. 5 (1987), pp. 742--747. Supplied attribution for the dimension-$2$ identity result when $p\neq 2$; no article title or link is supplied.
\item D. E.-C. Ben-Ezra and E. Zelmanov, \emph{Trans. Amer. Math. Soc.}, volume 374, no. 6 (2021), pp. 4093--4128. Supplied attribution for the dimension-$2$ pro-$2$ identity result in characteristic $2$; no article title or link is supplied.
\end{itemize}
\end{document}
